\textbf{Пример}

Рассмотрим следующий вызов:

\texttt{init(3, [2, 6, 9], [3, 1, 2], [2, 2, 3], [1, 0, 1])}

\includegraphics[scale=0.6]{1.png}

Диаграмма выше соответствует этому вызову. Каждый квадрат обозначает подземелье. Для подземелий $0$, $1$ и $2$, значения $s[i]$ и $p[i]$ записаны внутри квадратов.
Пурпурные стрелки показывают, куда герой перемещается в случае победы, а чёрные~--- куда герой перемещается в случае поражения.

Пусть грейдер вызвал функцию \texttt{simulate(0, 1)}.

Игра пройдет следующим образом:

\begin{center}
\renewcommand{\arraystretch}{1.5}
\begin{tabular}{|c|c|c|}
\hline
 Подземелье & Сила героя до сражения & Результат \\
\hline
 $0$ & $1$ & Поражение \\
\hline
 $1$ & $4$ & Поражение \\
\hline
 $0$ & $5$ & Победа \\
\hline
 $2$ & $7$ & Поражение \\
\hline
 $1$ & $9$ & Победа \\
\hline
 $2$ & $15$ & Победа \\
\hline
 $3$ & $24$ & Конец игры \\
\hline
\end{tabular}
\end{center}

Таким образом, функция должна вернуть $24$.

Пусть грейдер вызвал функцию \texttt{simulate(2, 3)}.

Игра пройдет следующим образом:

\begin{center}
\renewcommand{\arraystretch}{1.5}
\begin{tabular}{|c|c|c|}
\hline
 Подземелье & Сила героя до сражения & Результат \\
\hline
 $2$ & $3$ & Поражение \\
\hline
 $1$ & $5$ & Поражение \\
\hline
 $0$ & $6$ & Победа \\
\hline
 $2$ & $8$ & Поражение \\
\hline
 $1$ & $10$ & Победа \\
\hline
 $2$ & $16$ & Победа \\
\hline
 $3$ & $25$ & Конец игры \\
\hline
\end{tabular}
\end{center}

Таким образом, функция должна вернуть $25$.